\section{統計学的にデータを調べるということ}

  \subsection{わずかな差をしっかりみる}

    あなたはじゃんけんが強いだろうか.

    実はじゃんけんはランダムではない.
    ある調査\footnote{%
    芳沢光雄（桜美林大学）による, 学生725人・延べ11{,}567回のじゃんけんデータに基づく.
    新聞報道を通じて広く知られている.}
    によれば, グー・チョキ・パーの出現比は
    およそ 35 : 32 : 33. ほぼ互角に見えるが, 完全に一様ではない.
    そしてこの程度の差でも, データ数が十分に大きければ「偶然ではない」と言い切れる.
    つまり, 現実に勝つ確率が高いといえるのだ. 
    じゃんけんに統計的な正解があるとすれば, それはパーだ.

    さらに, あいこが絡む場合はもっと有効な手がある. この研究では, あいこの後に手を変える確率はおよそ 78\% とされているのだ.
    グーであいこになったなら, 相手の次の手はパーかチョキが高確率だ. グーはほぼ出てこない.
    ということはパーに勝つチョキを出せば, 負ける確率をかなり減らすことができる. 
    相手の手が読めるなら, それに合わせて自分の手を選べばよいのだ.
  
  \subsection{確率論と統計学}
    確率論と統計学は, 考える向きが反対である.

    確率論は, 不確かさそのものをあつかうための数学であり, 統計学は, 不確かさを前提に判断を下すための数学である.

    たとえばサイコロを振るとき, 次の目が何になるかは振ってみるまで分からない. 
    しかし「6の目が出る確率は1/6」「100回振ればだいたい何回6が出るか」といったことは, 
    振る前から計算できる. 
    このように不確かさの仕組みを理論的に組み立て, 
    そこから出てくるデータの振る舞いを予測するのが, 確率論で考えていることだ.

    一方で統計学では, 手元にデータはあるが, それを生んだ仕組みを直接調べるのは難しい, という状況を想定する. 
    たとえば100人にアンケートをとった結果から, 
    日本人全体の意見がどうなっているかを推し量ったり, 
    あるいはじゃんけんを11{,}567回観察した結果から, 
    人がじゃんけんで出す手にどんなクセがあるかを読み取ったりする. 
    データは出発点であり, これらは知りたい全体の情報からみると不完全なものだが, 
    それでもこのデータから背後の仕組みに対して何が言えるのか, ということを考えていくのが統計学である.

  \subsection{目標をかたちにする}

    とはいえ, 実際にデータを目の前にしてみると話は別である. 
    内心では「面倒だ」「とりあえず誰かに任せたい」と感じてしまうのが人情である. 
    しかし「これが分かれば儲かる仕事の判断材料になる」「ここに違和感の正体があるかもしれない」といった見通しが立つと, 
    ちょっとやってみるかな, という気になってくる（かもしれない）. 
    その分かれ目は, 「自分は何を知りたいのか」を言葉にできているかどうかにある.

    目的は最初から明確でなくてもいい. 「売上が伸びている気がする, なぜだろう」くらいの
    ぼんやりした問いで十分だ. 数値や図を見ながら問いを絞り込んでいく, その過程そのものが分析である.

    そもそも統計的分析に, 必ずこうしなければならないという手順はない. 
    データと目的次第で, 手法も順序も変わるし, 行きつ戻りつが普通だ. 
    本書では問いをどう立てるかという問題設定の方法論そのものには踏み込まないが, 
    どの章を読んでいるときも 「自分はいま何を明らかにしようとしているのか」を頭の片隅に置いておいてほしい. 
    それだけで, 手法の意味がずっと掴みやすくなるはずだ.

  \subsection{理解してほしいこと}

    本書では, データからどうやって有用な情報をとりだすか, という統計学の実学的な側面を扱う.
    このとき鍵となる概念はそれほど多くない. 
    それは平均と分散, そして正規分布である.
    この三つの概念がどういうものなのか, なぜ重要なのか, 
    それをつかめば統計学の入門としては概ね十分である.

  \subsection{分析の準備}

    本書を通じて, ひとつの共通したデータを手がかりに分析を進めていく.
    扱うのは「コンテンツの人気と評価」に関するデータだ.
    具体的な題材と変数の設計については後で紹介する.

    % ランニング例（題材・表・問い）をここに挿入

    ここで, 統計で使う最も基本的な言葉を確認しておこう.
    手元にあるデータは, 世界のごく一部を切り取ったものにすぎない.
    この切り取られた一部分であるデータを\textbf{標本}と呼び, その背後にある全体を\textbf{母集団}と呼ぶ.
    厄介なことに, 標本はどんな取り方をしても必ずどこかに偏りや情報の欠けを持ってしまう. 
    基本的に完全なデータなど存在しないのだ.
    統計学はそうした不完全な標本を出発点とする. だから,
    「どんなことを, どこまで確からしく語れるのか」についてはいつも意識しておかなければならない.

    各章では, その手法を説明するのに適したデータを使う. 章末ではこの例に戻ってきて, 学んだ手法を当てはめてみる. 全章を通した一貫した分析というには少し届かないかもしれないが, 同じデータを別の手法で見ることで, 各手法の役割がはっきりするはずだ.

\begin{tcolorbox}[
  colback=gray!5,
  colframe=gray!50,
  arc=4mm,
  boxrule=0.5pt,
  title={\textbf{コラム：データとは何か}},
  fonttitle=\large
]

「3, 5, 2, 6」という数字の並びがある.
さて, これは何だろう.
暗証番号かもしれないし, サイコロの出目かもしれない.
実はこれだけでは, 何もわからない.
「あるカフェの月曜から木曜までの売上（万円）」だと聞いて初めて,
この数字は意味のある情報になる.
データとは, 何について, どのように測ったかという文脈があって初めて成り立つものだ.

そして, 文脈があるということは, データにも性格の違いがあるということだ.
売上額や気温のように足し引きに意味がある\textbf{数値データ}もあれば,
天気（晴れ・雨・曇り）や血液型のように分類を表す\textbf{カテゴリデータ}もある.
ちょっと厄介なのは, 数字で書いてあるからといって数値データとは限らないことだ.
たとえばゲームのランキング順位は立派な数字だが,
1位と2位の差と2位と3位の差が同じ意味を持つとは限らない.
あれは量ではなく, 順序を示しているだけである.

性格だけでなく, 姿もさまざまだ.
画像, 音声, テキスト, 人と人のつながり.
データは表の形をしているとは限らない.
しかし本書で扱うのは, 行と列が並んだ, いわゆる表のデータである.
当たり前のようだが, 実はこれ自体がデータの世界のほんの一部にすぎない.
ただ, この「表のデータ」こそが統計の出発点であり,
本書で手を動かしながら付き合っていく相手だ.

\end{tcolorbox}

  % \subsection{本書での前提}

  %   さて, 本書では次のような状況を想定する.

  %   あなたはECサイトで商品を選ぶ際, 他の購入者のレビューや星評価を参考にするだろう.  
  %   「星評価が高い商品はよく売れる」「レビューが多い商品は信頼できる」といった感覚は, 多くの人が持っているかもしれない.  
  %   しかし, 本当にレビューの数や評価の高さは, 商品の売上と統計的に意味のある関係があるのだろうか.

  %   例えば, あなたが運営するECサイトで, 最近販売された商品 20 点について, 次のデータを記録したとする.

  %   \begin{tabular}{ccccc}
  %   \hline
  %   商品ID & 商品カテゴリー & レビュー数 & 星評価平均 & 購買数 \\
  %   \hline
  %   P001 & 家電 & 15 & 4.2 & 120 \\
  %   P002 & 書籍 & 5 & 3.8 & 30 \\
  %   P003 & ファッション & 25 & 4.5 & 200 \\
  %   P004 & 食品 & 10 & 4.0 & 80 \\
  %   P005 & 家電 & 8 & 3.5 & 50 \\
  %   P006 & 書籍 & 2 & 3.0 & 15 \\
  %   P007 & ファッション & 18 & 4.3 & 150 \\
  %   P008 & 食品 & 12 & 4.1 & 90 \\
  %   P009 & 家電 & 30 & 4.8 & 250 \\
  %   P010 & 書籍 & 7 & 3.9 & 40 \\
  %   P011 & ファッション & 22 & 4.4 & 180 \\
  %   P012 & 食品 & 9 & 3.7 & 60 \\
  %   P013 & 家電 & 10 & 3.9 & 70 \\
  %   P014 & 書籍 & 3 & 3.2 & 20 \\
  %   P015 & ファッション & 28 & 4.6 & 220 \\
  %   P016 & 食品 & 14 & 4.2 & 100 \\
  %   P017 & 家電 & 20 & 4.3 & 160 \\
  %   P018 & 書籍 & 6 & 3.6 & 35 \\
  %   P019 & ファッション & 16 & 4.1 & 130 \\
  %   P020 & 食品 & 11 & 4.0 & 75 \\
  %   \hline
  %   \end{tabular}

  %   一見, レビュー数や星評価平均が高い商品ほど購買数が多いように見える.  
  %   だが, その印象が統計的に意味のある関係を示しているのか, それとも偶然の変動で説明できる範囲なのかは, ここからが本題である.  
  %   本書では, この小さな表から実際に手を動かして, 「要約 → 形 → 標本分布 → 推定 → 検定 → 関係 → モデル」という道具立てを順に試し,  
  %   \textbf{売上を伸ばすための具体的な示唆}を, 数と図と文章で導いていく.

% 目的
% ECサイトの購買データを用いて, 商品の売上を伸ばすための具体的な示唆を得る.
%
% （注）統計学は規則性や構造を明らかにし, 必要に応じて因果関係の仮説を検討するための枠組みである.


















% もう少し実用的な例を挙げよう．
% 目の前に大きな湖があるとする．
% もしこの湖にどのくらい魚がいるか見積もってください，と言われたらどのような方法を考えるだろうか．
% 前節の推定をうまく活用すると次のような方法が考えられる．
% まず，魚を一定の数，たとえば100匹ほど捕まえる．
% 次にこの魚にタグをつけて一旦放流し，しばらく待つ．
% そして魚が十分自由に動き回ったと思ったら（一週間程度だろうか）また魚を先ほどと同じ数，100匹捕まえて，タグが付いた魚の数を数える．
% 数えてみた結果，その中にタグの付いた魚が10匹いたとしよう．割合は1/10である．
% するとこの湖には最初にタグをつけた100匹が1/10くらいの割合で存在する, とある程度の確からしさで推測できるので，魚全部ではおよそ1000匹ほどいるだろうと見積もることができる．

% \subsubsection{統計学の起源と進化 - 国家の数字から万人の道具へ}
% 「統計（statistics）」という言葉は, もともと「国家（state）の数字」を意味していたらしい．
% 17世紀のヨーロッパでは人口や税収, 疫病による死亡記録などが集められ始め, 人々はそのデータの中に隠れたパターンや法則があることに気づき始めた. 
% 1662年, ジョン・グラントはロンドンの死亡記録を整理し, 疫病の流行周期や平均寿命を初めて統計的に明らかにした. 
% これが近代統計学のはじまりとされている.

% 18世紀から19世紀にかけて, ラプラス, ガウス, ポアソンといった数学者たちによって確率論が発展し, 統計学の理論的基盤が形成された．
% 20世紀には, R.A.フィッシャーやカール・ピアソンらによって実験計画法や仮説検定の手法が確立され, 統計学は科学研究の必須ツールとなった．

% さらに看護師フローレンス・ナイチンゲールはクリミア戦争の現場で衛生状態を統計データにまとめ, 「ナイチンゲールのバラ」と呼ばれる図を用いて衛生改善の必要性を説いた. 
% 実はこの時代の戦争において, 戦闘行為による直接的な死者よりも, 衛生状況が悪いことによる負傷兵の死者の方が多かったのである. 
% これにより統計は単なる数字遊びではなく, 人々を説得し社会を動かす力を持つことが明らかになった. 

% 現代では, 統計学はビッグデータ, 機械学習, 人工知能といった先端分野を支える基盤となっている．
% かつては専門家だけのものだった高度な統計手法も, コンピュータの発達により広く一般に利用可能になり, 
% ビジネス, 医療, 教育, スポーツ, 芸術など, あらゆる領域で活用されている．

% \subsubsection{標本と統計}
% 改めて統計の言葉遣いについて改めて言及しておこう．
% データとは, 観察や測定によって得られた情報のことを指す. 
% データには様々な特性や属性があり, これらを「変数」と呼ぶ. 
% 変数には, 身長や体重のような数値で表される量的変数と, 性別や血液型のようなカテゴリーで表される質的変数があるが，これらについてはあとで詳しく述べる.
% そして，今，手元にある「データの集まり」のことを標本という．
% 標本は，実際は巨大な全体集合からその一部を取り出したものである．
% この全体集合のことを「母集団」と言い，母集団から標本を取り出すことを「抽出する」という．
% 標本の大きさ, つまりデータの数のことを「標本サイズ」または「サンプルサイズ」と呼ぶ. 例えば, 100人分のデータがあれば, サンプルサイズは100となる.
% 統計的な量（たとえば平均など）を計算したり, データの特徴を図や表で表したりして「今手元にあるデータの理解を深める」ことを「記述統計」と呼ぶ. 
% 一方, 標本から母集団の特徴を推測したり，データの比較をしたりすることを「推測統計」という.
% 統計の項目ではこれら記述統計と推測統計の基本的なこと，つまり何がしたくて，どうやって計算して，なにがわかるのかを理解することを目指す．



% \subsection{データを「見る」ための基礎 - 無秩序から秩序への変換}

% \textbf{数字の羅列に意味を付与する - コンテキストの魔法}\newline
% データは単に文字や数字が並んでいるものではない.
% その並びには必ずなんらかの背景や文脈が存在する.  
% 例えば, 「20, 50, 30, 90, 40, 80, 90」だけでは何を表しているのか意味不明である. 
% しかし, 「1週間の来店人数（月～日）：20, 50, 30, 90, 40, 80, 90人」と文脈が与えられると, 突然意味が生まれる. 
% 水曜日（30人）よりも火曜日（50人）の方が来店者が多く, 週末（土日：80, 90人）に集中していることが分かる.
% これにより, たとえば「週末は人が多いよね」となんとなく思っていたことが事実として現れてくるのである  

% このように, 数字に「タイトル（何についてのデータか）」「ラベル（各値が何を表すか）」「単位（何を数えているのか）」「時間的・空間的文脈（いつ, どこのデータか）」を付与することで, 無秩序な数字の羅列は, 私たちが「見て」解釈できる情報へと変換されるのです. 

% \textbf{「見えない」から「見える」へ：データの文脈化の技術}\newline
% 記録されたままの状態のデータを「生」と表現することがあるが, そのようなデータは非常に扱いにくい状態であることが多い. 
% それを磨き上げ, 形を整えることで初めて宝石としての価値が現れます. データを「見える」ようにするためには, 以下のような文脈化の技術が重要です：

% 1) \textbf{比較対象の設定}：東京の平均気温25℃という数字は, それだけでは「暑い」のか「寒い」のか判断できませんが, 過去10年の同時期の平均が22℃だと分かれば「例年より暑い」と理解できます. 同様に, ある美術展の来場者数5,000人も, 類似の展覧会の平均来場者数と比較することで, その「多さ」や「少なさ」が初めて意味を持ちます. 

% 2) \textbf{時間軸の導入}：データに時間的文脈を加えることで, 変化やトレンドが見えてきます. 例えば, 美術館の月間来場者数を単月で見るより, 過去5年間の推移として見ることで, 季節変動や長期的な人気の変化が明らかになります. 

% 3) \textbf{空間的・地理的文脈}：同じ「美術館の来場者数」でも, 都心に位置する国立美術館と地方の小規模美術館では, 比較の意味が異なります. 地理的条件や施設の規模などの空間的文脈を考慮することが, 公平な比較には不可欠です. 

% 4) \textbf{条件を追加する}：単一の変数だけでなく, 関連する複数の変数と組み合わせることで, より豊かな文脈が生まれます. 例えば, 美術館の来場者数と天気, 交通アクセス, 特別展の有無などを併せて分析することで, 来場者数の変動要因が見えてきます. 


% \textbf{データの尺度と比較可能性 - 同じ土俵で考える}\newline
% データを扱うとき, 必ず意識しなければならないのはそのデータの「単位」である. 
% 同じ「3」という数字でも「3人」と「3cm」と「3秒」では全く異なる情報となり, 比較したりすることがようやく可能となるのである. 
% また, さらに「3点」という評価も, 5点満点なのか10点満点なのかによって意味が変わってくる. 

% 統計学では, データの尺度を以下のように分類します：

% 1) \textbf{名義尺度（nominal scale）}：カテゴリーや分類を表す尺度（例：展示作品のジャンル, 来場者の性別）. 順序や大小関係はなく, 単なる区別のみ. 

% 2) \textbf{順序尺度（ordinal scale）}：順序や優劣を表す尺度（例：展示の満足度評価「とても良い, 良い, 普通, 悪い, とても悪い」）. 大小関係はあるが, 間隔の等間隔性は保証されない. 

% 3) \textbf{間隔尺度（interval scale）}：等間隔性を持つ尺度（例：温度（摂氏）, 日付）. 「ゼロ」という数字の意味がある程度自由に設定され, 絶対的な「無」を表さない. たとえば「摂氏 0 度」は温度が全くないということではないし, 4 月 0 日という日は存在しない.  

% 4) \textbf{比率尺度（ratio scale）}：絶対的なゼロ点を持つ尺度（例：身長, 重さ, 人数）. 比率の比較が意味を持つ（「AはBの2倍」など）. 

% 1, 2 のようなデータはデータどうしで「四則演算」ができないことに注意が必要である. 
% たとえば血液型「A型」と「B型」を足すことは一般的に見て意味不明だと思う. 
% このような場合はたとえば血液型をラベルのように用いて「A型」の「人数」などを数えることにより, 足すことに意味が生まれる. 

% 1,2 のようなデータはカテゴリデータ（あるいは質的データ）と呼ばれ, 3, 4のようなデータは数値データ（あるいは量的データ）と呼ばれる. 
% 本書では量的データの扱いを通して統計学の基礎を身に着けていこう. 
% 質的データに対する基本的な分析についても 3 節で触れる. 


% \subsection{データと表・視覚的整理 - 情報の構造化と可視化}

% \textbf{表が果たす強力な役割 - 情報の地図作り}\newline
% データというと以下のような表を思い浮かべる方も多いのではないだろうか. 
% 大げさに聞こえるかもしれないが, 表は情報を整理するための非常に優れたツールである. 
% 列の名前を見るだけでだいたいどんなことを表すデータなのか把握でき, ピンポイントでほしい情報をすぐに見つけることができる. 

% \begin{tabular}{|l|c|c|c|c|c|}
% \hline
% & 午前 & 昼食時 & 午後 & 夕方 & 合計 \\
% \hline
% 月曜日 & 15 & 20 & 25 & 10 & 70 \\
% 火曜日 & 35 & 45 & 50 & 30 & 160 \\
% 水曜日 & 30 & 40 & 45 & 25 & 140 \\
% 木曜日 & 25 & 35 & 40 & 20 & 120 \\
% 金曜日 & 40 & 50 & 55 & 35 & 180 \\
% 土曜日 & 80 & 100 & 120 & 60 & 360 \\
% 日曜日 & 90 & 110 & 130 & 70 & 400 \\
% \hline
% 合計 & 315 & 400 & 465 & 250 & 1430 \\
% \hline
% \end{tabular}

% この表からは, 「週末（土日）の来場者数が平日の約3倍である」「全体的に午後の時間帯が最も混雑する」「月曜日は最も来場者が少なく, その差は日曜日の約1/6である」といった情報が一目で読み取れます. このような表の「読み方」を身につけることは, データを「見る」ための基本的なスキルです. 

% \textbf{整理の技術：データの見せ方が結論を変える}\newline  
% 表の作り方（ランキング順, 時系列順, カテゴリー別など）の工夫により, データの読み取り方や印象は大きく変わります. 例えば, 同じ美術館の来場者データでも, 「曜日別の集計」と「天気別の集計」では, まったく異なる傾向が浮かび上がります. 

% さらに, どの軸で並べるか, どの値を強調するかによっても印象は変化します. 例えば上記の表を「来場者数の多い順」に並べ替えると：

% \begin{tabular}{|l|c|c|c|c|c|}
% \hline
% & 午前 & 昼食時 & 午後 & 夕方 & 合計 \\
% \hline
% 日曜日 & 90 & 110 & 130 & 70 & 400 \\
% 土曜日 & 80 & 100 & 120 & 60 & 360 \\
% 金曜日 & 40 & 50 & 55 & 35 & 180 \\
% 火曜日 & 35 & 45 & 50 & 30 & 160 \\
% 水曜日 & 30 & 40 & 45 & 25 & 140 \\
% 木曜日 & 25 & 35 & 40 & 20 & 120 \\
% 月曜日 & 15 & 20 & 25 & 10 & 70 \\
% \hline
% 合計 & 315 & 400 & 465 & 250 & 1430 \\
% \hline
% \end{tabular}

% このように並べ替えると, 「曜日による来場者数の差が非常に大きい」ことがより明確になります. 
% データの整理方法自体が, 分析の視点と結論を左右する重要な要素なのです. 

% 統計学では, このような「データの見せ方」そのものが分析の一部であることを認識し, 目的に応じた適切な整理方法を選択することが重要です. 
% 異なる整理方法で同じデータを複数回見ることで, 多角的な理解が可能になります. 

% ほかにも実はスマートフォンに保存されている写真などの画像データは「表」そのものである. 
% 各セルの大きさは「ピクセル」と呼ばれ, それぞれのセルにはそのセルの「色情報」が保存されている. 
% これを特殊な読み取り方をすることによって画像として表示されるのである. 

% \textbf{表からグラフへ：視覚化の第一歩}\newline
% 整理された表は, 数字に意味を与え, 後のグラフや図表による可視化の基盤となります. 
% 人間の脳は, 数字の羅列よりも視覚的なパターンを認識するのが得意です. 
% 例えば, 「月別売上データ」は表よりも折れ線グラフで示した方が, 季節変動や成長トレンドが一目で理解できます. 

% グラフは表の情報を視覚言語に翻訳するものであり, その選択は伝えたいメッセージによって異なります：

% 1) \textbf{時系列データの変化}：折れ線グラフが最適（例：年間を通じた来場者数の変動）

% 2) \textbf{部分と全体の関係}：円グラフや積み上げ棒グラフが有効（例：来場者の年齢層構成）

% 3) \textbf{分布の形状}：ヒストグラムや箱ひげ図が適切（例：滞在時間の分布）

% 4) \textbf{相関関係}：散布図が直観的（例：入場料と来場者数の関係）

% この「表からグラフへ」の変換プロセスは, データを「見る」技術の核心部分です. 
% 適切なグラフ形式を選択し, タイトル, ラベル, 凡例, スケールなどを工夫することで, データの持つメッセージをより効果的に伝えることができます. 

% \textbf{デジタル時代の表現技術 - データ体験の革新}\newline
% 現代のデータ可視化技術は, 静的な表やグラフを超え, インタラクティブなダッシュボード, リアルタイム更新, 3D表現, 拡張・仮想現実（AR/VR）を用いた没入型データ体験など, 多様な方法でデータを「体験」できるツールへと進化しています. 

% 例えば, 美術館の来場データを分析する場合, 従来の静的なグラフだけでなく, 以下のような革新的な可視化が可能になっています：

% 1) \textbf{インタラクティブマップ}：館内の平面図上に来場者の動線や滞在時間をヒートマップで表示し, 人気の展示や混雑エリアを直感的に把握

% 2) \textbf{時間変化アニメーション}：一日の時間帯による来場者の流れを動画のように表現し, 混雑の波を視覚的に理解

% 3) \textbf{多次元データの可視化}：来場者の年齢, 性別, 滞在時間, 支出金額などの多変量データを, 色, 形, 大きさなどの視覚要素に変換して複合的に表現

% これらの先進的なデータ表現技術を効果的に活用するためには, 統計学の基本原則を理解することが不可欠です. 技術は進化しても, 「データを意味ある形で整理し, パターンを見出し, 解釈する」という統計学の本質は変わりません. 現代の視覚化ツールは, この本質的なプロセスをより効率的かつ魅力的に実行するための道具なのです. 

% 統計学の基礎を身につけることは, 現代のデータ豊富な環境で「見る力」を養うための重要なステップです. 
% 数字の羅列から意味のあるストーリーを引き出す能力は, ビジネス, 研究, 芸術, そして日常生活のあらゆる場面で, より良い判断と創造的な問題解決につながるでしょう. 




